% Retained original regret lower bounds, for insertion into the main source.
% Original source: mab_aa_main.tex lines 1002--1179 and 1215--1355.
% The local-alternative proof is replaced; no original UCB-AA upper claim is used.
\section{Regret lower bounds in representative arrival regimes}
\label{sec:retained_regret_lower_bounds}

We retain two lower-bound regimes from the original regret formulation.
The first describes the evidence needed to reject inferior arrivals when
the best arm is already available. Its static special case is a classical
calibration \citep{LAI19854,BurnetasKatehakis1996}. The second describes
successive local comparisons as the benchmark improves. Both concern regret
at a deterministic physical horizon. They are separate from the reporting
requirements introduced below, and are not claimed as matching converses to
an upper bound evaluated at a random round-completion time.

\subsection{Static arms and arriving distractors}

For each horizon $T$, fix an exogenous deterministic arrival schedule and a
baseline environment $\nu^{(0)}$. Arm $0$ is available from time $1$ and has
mean $r^\star$ strictly larger than every other mean. Let $\mathcal I_T$ be
a nonempty collection of inferior arrivals. For $i\in\mathcal I_T$, write
\[
 r_i=r^\star-\Delta_i,\qquad
 H_i(T)=T-\tau_i+1,\qquad
 H_{\min}(T)=\min_{i\in\mathcal I_T}H_i(T).
\]
The reward law of arm $i$ is $P_i$. Suppose there is a law $Q_i$ on $[0,1]$
with mean $r^\star+\eta_i$, where $\eta_i>0$, and put
$d_i=D(P_i\Vert Q_i)$. Let $\nu^{(i)}$ change only this arm's law to $Q_i$,
retaining all other reward laws and the arrival schedule. Set
$\mathcal V_T=\{\nu^{(0)}\}\cup\{\nu^{(i)}:i\in\mathcal I_T\}$.

\begin{theorem}[Local-alternative regret lower bound]
\label{thm:retained_distractors}
Assume $H_{\min}(T)\to\infty$ and that fixed positive constants
$\underline\Delta,\underline\eta,\overline d$ satisfy, for all sufficiently
large $T$ and every $i\in\mathcal I_T$,
\[
 \Delta_i\ge\underline\Delta,\qquad
 \eta_i\ge\underline\eta,\qquad
 0<d_i\le\overline d.
\]
Suppose a sequence of nonanticipating policies $\{\pi^{(T)}\}$ is uniformly
good on this family at scale $H_{\min}(T)$, meaning
\begin{equation}
 \sup_{\nu\in\mathcal V_T}\E_\nu
       R_T(\pi^{(T)};\nu)=o\bigl(H_{\min}(T)^a\bigr)
       \quad\text{for every }a>0.
 \label{eq:retained_uniform_goodness}
\end{equation}
Then, for all sufficiently large $T$,
\begin{equation}
 \E_{\nu^{(0)}}R_T(\pi^{(T)};\nu^{(0)})
 \ge \frac1{12}\sum_{i\in\mathcal I_T}
                 \frac{\Delta_i}{d_i}\log H_i(T)
 \ge \frac{\underline\Delta}{12\overline d}
            |\mathcal I_T|\log H_{\min}(T).
 \label{eq:retained_distractor_lb}
\end{equation}
\end{theorem}

The conclusion is conditional on \eqref{eq:retained_uniform_goodness}; that
condition can be restrictive for growing families. Indeed, the theorem
requires $|\mathcal I_T|\log H_{\min}(T)=o(H_{\min}(T)^a)$ for every $a>0$
if such a policy sequence exists. It does not assert existence in arrival
regimes that violate this necessary compatibility condition. Equally spaced
arrivals with spacing $\lfloor T/|\mathcal I_T|\rfloor\to\infty$ give
$H_{\min}(T)\ge\lfloor T/|\mathcal I_T|\rfloor$. If every arrival occurs
before $\alpha T$ for a fixed $\alpha<1$, then
$H_{\min}(T)\ge(1-\alpha)T$. These substitutions describe the corresponding
lower bounds without removing the uniform-goodness condition.

For a fixed collection of $K$ inferior arms all available at time $1$, take
$H_i(T)=T$. If the local alternatives satisfy
$d_i\le C_i\Delta_i^2$ with fixed $0<C_i<\infty$, then the same argument gives
\begin{equation}
 \E_{\nu^{(0)}}R_T(\pi^{(T)};\nu^{(0)})
 \ge \frac1{12}\sum_{i=1}^K\frac{\log T}{C_i\Delta_i}
 \label{eq:retained_static_lb}
\end{equation}
for all sufficiently large $T$. The laws, gaps, and $K$ are fixed in this
specialization. These are applications of the usual bandit change-of-measure
argument \citep{kaufmann2016}, not new claims about a regret comparator.
\suppproof

\subsection{Successive near-tied benchmark improvements}

Fix integers $J\ge1$ and $H_1,\ldots,H_J\ge1$ with
$\sum_{m=1}^JH_m=T$. Write $s_0=0$,
$s_m=\sum_{\ell=1}^mH_\ell$, and
$\mathcal B_m=\{s_{m-1}+1,\ldots,s_m\}$. At the start of block $m$,
two Bernoulli arms $c_m,d_m$ arrive and remain available, with means
\[
 r_{c_m}=\theta_m,\qquad
 r_{d_m}=\theta_m+\sigma_m\varepsilon_m,
 \qquad \sigma_m\in\{-1,+1\}.
\]
There are no other arms in this construction, so $K_T=2J$. Fix
$0<\gamma<1/2$ and require
\begin{equation}
 \theta_m-\varepsilon_m,\ \theta_m+\varepsilon_m
       \in[\gamma,1-\gamma],\qquad \varepsilon_m>0,
 \qquad
 \theta_{m+1}\ge\theta_m+\varepsilon_m\quad(m<J).
 \label{eq:retained_progressive_parameters}
\end{equation}
Let $\mathcal V_T^{\rm prog}$ contain the $2^J$ sign configurations.
The local perturbations $\varepsilon_m$ in this construction are not the
near-optimality tolerance used later for reporting.

\begin{theorem}[Deterministic-block regret lower bound]
\label{thm:retained_progressive}
Every nonanticipating policy satisfies
\begin{equation}
 \sup_{\nu\in\mathcal V_T^{\rm prog}}\E_\nu R_T(\pi;\nu)
 \ge\frac12\sum_{m=1}^JH_m\varepsilon_m
       \left(1-\sqrt{\frac{2H_m\varepsilon_m^2}
                                {\gamma(1-\gamma)}}\right)_+.
 \label{eq:retained_progressive_lb}
\end{equation}
In particular, if $H_m\varepsilon_m^2\le\kappa\gamma(1-\gamma)$ for all
$m$, with $0<\kappa<1/2$, the right side is at least
$c_\kappa\sum_m H_m\varepsilon_m$, where
$c_\kappa=(1-\sqrt{2\kappa})/2>0$. The same bounds hold after taking the
infimum over policies on the left.
\end{theorem}

If $H_m\asymp T/J$, $\varepsilon_m\asymp\sqrt{J/T}$, and the constants
satisfy the preceding small-information condition, this gives a lower bound
of order $\sqrt{JT}$. Bounded rewards restrict this specialization. If
$\varepsilon_m\ge c_\varepsilon\sqrt{J/T}$ for every $m$, then
\[
 Jc_\varepsilon\sqrt{J/T}
 \le \sum_{m=1}^J\varepsilon_m
 \le \theta_J+\varepsilon_J-(\theta_1-\varepsilon_1)
 \le1-2\gamma.
\]
Consequently $J=O(T^{1/3})$ is necessary for this family. This is a
feasibility restriction of the bounded ladder construction, not an
impossibility boundary for the unrestricted arriving-arm model. No
uniform-goodness assumption is needed for
\eqref{eq:retained_progressive_lb}. \suppproof
